\lecture{10}
\title{Consistency of the MLE}

\begin{document}

\subsection{Kullback-Leibler (KL) Divergence}

Maximizing the likelihood turns out to be closely related to
minimizing the ``distance'' between
$\mathbb{P}_{\theta^*}$ and $\mathbb{P}_{\theta}$.
Here's how we formalize that notion of ``distance''.

\textbf{Definition.} Given distributions $\mathbb{P}_1$ and $\mathbb{P}_2$
with PDFs $f_1(x)$ and $f_2(x)$, their \textbf{Kullback-Leibler (KL) divergence} is:
\begin{align*} D_{\mathrm{KL}}(\mathbb{P}_1 \parallel \mathbb{P}_2) :=
\mathbb{E}_{f_1}\left[ \log\left(\frac{f_1(X)}{f_2(X)}\right)\right] & =
\int_{\mathcal{R}} f_1(x) \log \left(\frac{f_1(x)}{f_2(x)}\right) \, \mathrm{d}x \\ & =
\left( \int_{\mathcal{R}} f_1(x) \log f_1(x) \, \mathrm{d}x \right) - 
\left( \int_{\mathcal{R}} f_1(x) \log f_2(x) \, \mathrm{d}x \right). \end{align*}
Here, $\mathcal{R}$ is the entire support of $\mathbb{P}_1$,
which is also the entire support of $\mathbb{P}_2$.
Also, as usual, the $\mathbb{E}_{f_1}$ notation
indicates that the random variable $X$ is to be sampled using the PDF $f_1$.

Note that $D_{\mathrm{KL}}$ is not a metric; it is not symmetric, nor does it
satisfy the triangle inequality.  But Jensen's inequality implies it is at least nonnegative,
with $D_{\mathrm{KL}}(\mathbb{P}_1 \parallel \mathbb{P}_2) = 0$ if and only if
$\mathbb{P}_1 = \mathbb{P}_2$.

\begin{remark}
    Jensen's inequality states that
    for any concave function $\varphi$ and any random variable $Z$,
    we have $\mathbb{E}[\varphi(Z)] \leq \varphi(\mathbb{E}[Z])$.
    Equality is precisely when $Z$ is almost-surely constant,
    assuming $\varphi$ is \emph{strictly} concave.
\end{remark}

Importantly, this implies the following.

\begin{theorem}{KL Divergence Minimization}
    The map $\theta \mapsto
    D_{\mathrm{KL}}(\mathbb{P}_{\theta^*} \parallel \mathbb{P}_{\theta})$
    is minimized when $\theta = \theta^*$.
    Also, the population log likelihood
    $\ell(\theta) := \mathbb{E}_{\theta^*}[\log f_{\theta}(X)]$
    is maximized at $\theta = \theta^*$.
\end{theorem}

\begin{proof}
    The first sentence is true by Jensen's; the point is that $\log$ is strictly concave, so:
    \[ D_{\mathrm{KL}}(\mathbb{P}_{\theta^*} \parallel \mathbb{P}_{\theta})
    = -\mathbb{E}_{\theta^*} \left[ \log \left(\frac{f_{\theta}(X)}{f_{\theta^*}(X)}\right) \right]
    \geq -\log \left( \mathbb{E}_{\theta^*} \left[ \frac{f_{\theta}(X)}{f_{\theta^*}(X)} \right] \right)
    = -\log(1) = 0. \]
    Equality occurs precisely when $\frac{f_{\theta}(X)}{f_{\theta^*}(X)}$
    (with $X \sim \mathbb{P}_{\theta^*}$) is almost-surely constant,
    or when $f_{\theta} \cong f_{\theta^*}$, as desired.

    The second sentence follows because minimizing
    $D_{\mathrm{KL}}(\mathbb{P}_{\theta^*} \parallel \mathbb{P}_{\theta})$
    is equivalent to maximizing
    $\int_{\mathcal{R}} f_{\theta^*}(x) \log f_{\theta}(x) \, \mathrm{d}x$;
    see the far-RHS of the definition of $D_{\mathrm{KL}}$.
    But $\int_{\mathcal{R}} f_{\theta^*}(x) \log f_{\theta}(x) \, \mathrm{d}x = \ell(\theta)$
    identically. ~ $\blacksquare$
\end{proof}

\subsection{MLE Consistency}

The upshot of introducing KL Divergence is the second sentence of the
``KL Divergence Minimization'' theorem.
We will leverage this fact to prove that the MLE is a consistent estimator.

\begin{theorem}{MLE Consistency}
    Under mild regularity conditions (identifiability in particular),
    the MLE is consistent; that is, if $\hat{\theta}_n$ is the MLE of the parameter $\theta$,
    then $\hat{\theta}_n \convprob \theta^*$.
\end{theorem}

\begin{proof}
    Recall the following.
    \begin{itemize}
        \item The MLE $\hat{\theta}_n$ is
        the maximizer of the log likelihood $\ell_n(\theta)$ by definition.
        \item The true parameter value $\theta^*$ is
        the maximizer of the population log likelihood $\ell(\theta)$
        by the previous theorem.
        \item For \emph{any} value $\theta'$,
        by the LLN, the sequence $\{\frac{1}{n}\ell_n(\theta')\}_{n = 1}^{\infty}$
        converges in probability to the constant $\ell(\theta')$.
    \end{itemize}
    These three points together imply $\hat{\theta}_n \convprob \theta^*$.
    Roughly speaking, the third bullet point says $\frac{1}{n}\ell_n \convprob \ell$,
    so the maximizer of $\ell_n$ should also converge in probability
    to the maximizer of $\ell$.
    (Yes, this needs to be justified---handwaving a little here.) ~ $\blacksquare$
\end{proof}

\end{document}